\subsection{反不变元素}

我们沿用定理 3 的假设与记号，并假设 K 的特征为 0。称 S 的一个元素 z 是 G 下反不变的，如果

\[
g (z) = (\det g) ^ {- 1} z
\]

对所有 \(g\in G\) 成立

设 H 是属于 G 且异于 1 的伪反射所成的集合。对所有 \(g \in H\) ，存在 \(e_{g} \in V\) 与 \(f_{g} \in V^{*}\) 使得

\[
g (x) = x + f _ {g} (x) e _ {g} \quad \text {   for   all   } x \in V.
\]

命题 5. (i) 设 \(\mathrm{D}\) 是 \(\prod_{g \in \mathrm{H}} e_g\) （\(\mathrm{S}\) 的一个元素）。\(\mathrm{S}\) 中在 \(\mathrm{G}\) 下反不变的元素是 \(\mathrm{RD}\) 的元素。

\begin{enumerate}
\def\labelenumi{(\roman{enumi})}
\setcounter{enumi}{1}
\tightlist
\item
  把 S 与多项式代数 \(\mathrm{K}[\mathrm{X}_1,\dots ,\mathrm{X}_l]\) 等同（通过选取 V 的一个基 \((\mathrm{X}_1,\dots ,\mathrm{X}_l)\) ），并设 \((\mathrm{P}_1,\dots ,\mathrm{P}_l)\) 是 S 中生成代数 R 的代数无关齐次元素（定理 3）。则 Jacobi 行列式 \(\mathrm{J} = \det \left(\frac{\partial\mathrm{P}_t}{\partial\mathrm{X}_j}\right)\) 具有 \(\lambda D\) 的形式，其中 \(\lambda \in K^*\) 。
\end{enumerate}

\begin{enumerate}
\def\labelenumi{\alph{enumi})}
\tightlist
\item
  按照 (ii) 中的记号，我们有
\end{enumerate}

\[
d \mathrm{P} _ {1} \wedge d \mathrm{P} _ {2} \wedge \dots \wedge d \mathrm{P} _ {l} = \mathrm{J} d \mathrm{X} _ {1} \wedge d \mathrm{X} _ {2} \wedge \dots \wedge d \mathrm{X} _ {l},
\]

于是，对所有 \(g \in G\) ，

\[
\begin{array}{c} g (\mathrm{J}) (\det g)   d \mathrm{X} _ {1} \wedge \ldots \wedge d \mathrm{X} _ {l} = g (\mathrm{J})   d (g \mathrm{X} _ {1}) \wedge \ldots \wedge d (g \mathrm{X} _ {l}) \\ = g (d \mathrm{P} _ {1} \wedge \ldots \wedge d \mathrm{P} _ {l}) = d \mathrm{P} _ {1} \wedge \ldots \wedge d \mathrm{P} _ {l} = \mathrm{J}   d \mathrm{X} _ {1} \wedge \ldots \wedge d \mathrm{X} _ {l}, \end{array}
\]

从而 J 在 G 下反不变。另一方面，S 的分式域 N 是 R 的分式域 E 的一个 Galois 扩张（第 2 号）；若 \(\Delta\) 是 E 的一个取值于 N 的扩域 \(\Omega\) 中的求导，则 \(\Delta\) 开拓为 N 的一个取值于 \(\Omega\) 中的求导（《代数》第 V 章，§ 16，第 4 号，定理 3）；由于诸 \(P_{i}\) 代数无关，由此可知 \(dP_{1} \wedge \ldots \wedge dP_{l} \neq 0\) ，从而 \(J \neq 0\) 。

\begin{enumerate}
\def\labelenumi{\alph{enumi})}
\setcounter{enumi}{1}
\item
  设 \(z\) 是 \(S\) 的一个在 \(G\) 下反不变的元素。我们证明 \(z\) 被 \(D\) 整除（在 \(S\) 中）。设 \(a\) 是 \(V\) 中一个非零向量。\(G\) 中保持 \(Ka\) 稳定的元素保持一个补超平面 \(L\) 稳定（附录，命题 2）；\(G\) 的一个保持 \(Ka\) 稳定的元素是 1 或以 \(a\) 为向量的伪反射，当且仅当它在 \(L\) 上诱导 1；于是以 \(a\) 为向量且属于 \(G\) 的伪反射与 1 一起构成一个循环子群 \(G'\) （\(G\) 的）；设 \(t\) 是它的阶。存在一个基 \((X_1, \ldots, X_l)\) （\(V\) 的），使得 \(a = X_1\) ，\(X_2 \in L, \ldots, X_l \in L\) ，且 \(z\) 可与一个多项式 \(P(X_1, \ldots, X_l)\) （系数在 \(K\) 中）等同。由 \(g(z) = (\det g)^{-1}z\) （对 \(g \in G'\) ），我们看到 \(X_1\) 在 \(P(X_1, \ldots, X_l)\) 中只以模 \(t\) 同余于 -1 的指数出现。于是 \(P(X_1, \ldots, X_l)\) 被 \(X_1^{t-1} = a^{t-1}\) 整除。而现在 \(D\) 是诸 \(a^{t-1}\) （对这些 \(a \in V\) ，\(t > 1\) ）的乘积（至多差一个标量因子），且 \(S\) 的这些元素两两互素。由于 \(S\) 是唯一分解的，\(z\) 被 \(D\) 整除。
\item
  由 a) 与 b)，J 在 S 中被 D 整除。而现在
\end{enumerate}

\[
\deg \mathrm{J} = \sum_ {i = 1} ^ {l} (k _ {i} - 1) = \operatorname{Card} (\mathrm{H})
\]

（命题 3），故 \(\deg J = \deg D\) ，从而 \(J = \lambda D\) ，其中 \(\lambda \in K\) 。由于 \(J \neq 0\) ，\(\lambda \in K^*\) 。这就证明了 (ii)。

\begin{enumerate}
\def\labelenumi{\alph{enumi})}
\setcounter{enumi}{3}
\tightlist
\item
  证明的 a) 与 c) 部分表明 D 在 G 下反不变。其次，若 \(y \in R\) ，则显然 yD 在 G 下反不变。最后，若 \(z \in S\) 在 G 下反不变，我们在 b) 中已见存在 \(y \in S\) 使 z = yD。由于 S 是整的，\(y \in R\) 。这就证明了 (i)。
\end{enumerate}
% semantic-fragments: legacy-input-seam
\relax{}
